\section{Описание реализации}
\label{sec:impl}

\subsection{Структуры данных и модули}

Программа реализована на языке Python с использованием библиотек NumPy, Pillow и Matplotlib. Основные компоненты приведены в таблице~\ref{tab:modules}.

\begin{table}[H]
	\centering
	\caption{Основные компоненты реализации}
	\label{tab:modules}
	\begin{tabularx}{\textwidth}{@{}lX@{}}
		\hline
		Компонент & Назначение \\
		\hline
		\texttt{image\_processing.py} & загрузка BMP, нормировка, построение позиций блоков, извлечение векторов, сборка изображения с усреднением перекрытий, формулы $Z$ и $K$; \\
		\texttt{network.py} & класс \texttt{RecirculationNetwork}: прямое распространение, адаптивный шаг, обновление весов, обучение по эпохам; \\
		\texttt{main.py} & интерактивный ввод параметров, опорный запуск, сохранение BMP, весов, метрик и графика ошибки; \\
		\texttt{experiments.py} & серии A--D, сравнение перекрытия, сохранение JSON и графиков для отчёта; \\
		\texttt{generate\_samples.py} & генерация трёх воспроизводимых BMP-примеров. \\
		\hline
	\end{tabularx}
\end{table}

Основные структуры данных:
\begin{itemize}
	\item \texttt{img\_array}~--- массив формы $(h,w,c)$;
	\item \texttt{positions}~--- список координат $(y,x)$ левых верхних углов блоков;
	\item $X$~--- матрица нормированных блоков размера $Q \times n$;
	\item $W$, $W'$~--- матрицы кодера и декодера размеров $n \times p$ и $p \times n$;
	\item $Y$~--- матрица сжатых представлений $Q \times p$;
	\item \texttt{errors}~--- история суммарной ошибки по эпохам;
	\item \texttt{metrics.json}, \texttt{experiments.json}~--- параметры, ошибки, $Z$, $K$, времена и результаты серий.
\end{itemize}

\subsection{Схема алгоритма обучения}

Схема реализованного алгоритма приведена на рисунке~\ref{fig:algorithm}.

\begin{figure}[H]
	\centering
	\begin{tikzpicture}[
		node distance=7mm,
		box/.style={draw, rounded corners, align=center, text width=9.2cm, inner sep=4pt, font=\small},
		dec/.style={draw, align=center, text width=8.4cm, inner sep=4pt, font=\small},
		arr/.style={-{Latex[length=2mm]}, thick}
	]
		\node[box] (load) {Загрузка BMP, нормирование пикселей в $[-1,1]$};
		\node[box, below=of load] (blocks) {Формирование позиций блоков (полное покрытие) и матрицы $X$};
		\node[box, below=of blocks] (init) {Инициализация $W$ и $W'$ случайными значениями};
		\node[box, below=of init] (epoch) {Новая эпоха: перемешивание порядка блоков};
		\node[box, below=of epoch] (sample) {Для очередного блока: $y=xW$, $\hat{x}=yW'$, $e=\hat{x}-x$, $\alpha=\beta/(\|x\|^2+\|y\|^2)$};
		\node[box, below=of sample] (update) {Обновление $W' \leftarrow W'-\alpha y^{\mathsf{T}}e$, \; $W \leftarrow W-\alpha x^{\mathsf{T}}\delta$};
		\node[dec, below=of update] (all) {Все блоки эпохи обработаны?};
		\node[box, below=of all] (err) {Вычисление $E$ по всей выборке};
		\node[dec, below=of err] (stop) {$E \le \varepsilon$ или $T=T_{\max}$?};
		\node[box, below=of stop] (finish) {Восстановление изображения, расчёт MSE, PSNR, $Z$, $K$; сохранение результатов};

		\draw[arr] (load) -- (blocks);
		\draw[arr] (blocks) -- (init);
		\draw[arr] (init) -- (epoch);
		\draw[arr] (epoch) -- (sample);
		\draw[arr] (sample) -- (update);
		\draw[arr] (update) -- (all);
		\draw[arr] (all) -- (err);
		\draw[arr] (err) -- (stop);
		\draw[arr] (stop) -- (finish);

		\draw[arr] (all.east) -- ++(1.2,0) |- node[pos=0.25, right, font=\scriptsize] {нет} (sample.east);
		\draw[arr] (stop.east) -- ++(1.2,0) |- node[pos=0.25, right, font=\scriptsize] {нет} (epoch.east);
		\node[right=1.35cm of stop, font=\scriptsize, align=left] {да};
	\end{tikzpicture}
	\caption{Схема алгоритма обучения}
	\label{fig:algorithm}
\end{figure}

\subsection{Порядок выполнения}

Загрузка изображения и построение позиций блоков выполняются один раз на запуск. Внутри эпохи веса изменяются после каждого блока; ошибка проверяется только после полного прохода по выборке. При достижении порога или лимита эпох формируется восстановленное изображение, вычисляются метрики и сохраняются веса, коды и график ошибки.

Модуль \texttt{experiments.py} использует ту же реализацию сети и отличается только организацией серий: для каждого набора параметров создаётся новая сеть, обучение выполняется с фиксированным seed, а результаты агрегируются в \texttt{experiments.json}. Порог $\varepsilon$ для серий A--C фиксируется по траектории опорного запуска (значение ошибки на 10-й эпохе), чтобы сравнивать варианты при одинаковом требовании к каждому изображению.
